\section{Background: the definitions, in the textbook's order}\label{sec:background}
% claims used here: au-1 au-2 C1 C2 C3 C4 C5 C12 C13 C14 C15 C16 C17 C18 C19 C20 C21 C22 C23 C30 C31 C32 C33
% C24 C25 m1-1 m1-2 m2-1 m2-2 m2-3 m2-4 m2-5 m2-6 m2-9 m2-10 m2-11 m2-12 m2-13 m2-14 m2-15 m3-1 m3-2 m3-3

Let $G=(V,E)$ be a finite simple undirected graph with $N$ vertices
$n_1,\dots,n_N$; let $u,v,w$ be vertices and $e$ an edge. These are the letters of
\Cref{tab:notation}, and every definition below uses exactly them. The
definitions follow the order of the course textbook~\cite{petrosyan2015grafneri}:
graph, adjacency, adjacency matrix, degree, walk, closed walk, trail, path, cycle,
simple cycle, distance, and then connectedness. Each definition cites the
textbook's number and printed page. After each definition come a remark on what it
allows and forbids, an example, and a non-example.

The three answers need these definitions: graph, adjacent, adjacency matrix and degree (Answer~1 and
the proof of the test), walk and trail and path (all three answers), closed walk (Answer~1), connected
and component (Answer~2), and the walk-to-path lemma. \emph{Cycle}, \emph{simple cycle} and \emph{distance}
are context: the answers use them only in remarks, examples and one by-product (\Cref{cor:components}).

Four graphs supply the examples. The first is the graph $B$ of the board matrix:
vertices $1,2,3$ (meaning $n_1,n_2,n_3$) and edges $12$ and $23$
(\Cref{fig:village}, left). The second is the board drawing $T$, a triangle with
edges $12$, $13$, $23$. The third is a small graph of this paper's own, $H$: a triangle $1,2,3$ with a
fourth vertex attached to vertex $3$, so its edges are $12$, $13$, $23$ and $34$ (\Cref{fig:H}). The fourth is the textbook's own example graph, which
this paper calls $G_{\mathrm{tb}}$ (\Cref{fig:textbook}); the seven edges used from it are $v_1v_2$, $v_1v_3$, $v_2v_3$, $v_2v_4$, $v_3v_4$, $v_3v_5$ and $v_4v_5$, the ones that the textbook's own examples use. $B$ has no cycle, and in $B$ every trail is a path (\Cref{rem:Bsimple}), so $B$ cannot show
the differences between the later definitions; $H$ and $G_{\mathrm{tb}}$ can.

\begin{figure}[ht]
\centering
% The graph H: a triangle 1,2,3 with a fourth vertex joined to vertex 3.
\setlength{\unitlength}{1pt}
\begin{picture}(100,55)(-50,-8)
  \put(0,0){\circle*{6}} \put(40,0){\circle*{6}} \put(0,40){\circle*{6}} \put(-40,40){\circle*{6}}
  \put(0,0){\line(1,0){40}}      % 1 2
  \put(0,0){\line(0,1){40}}      % 1 3
  \put(0,40){\line(1,-1){40}}    % 3 2
  \put(-40,40){\line(1,0){40}}   % 4 3
  \put(-3,-12){$1$} \put(37,-12){$2$} \put(3,44){$3$} \put(-43,44){$4$}
\end{picture}

\caption{The graph $H$ defined by this paper: a triangle on the vertices $1,2,3$ and a fourth vertex $4$ joined to $3$.}
\label{fig:H}
\end{figure}

\begin{figure}[ht]
\centering
% The seven edges forced by the textbook's own examples (a subgraph of its Fig. 1.2.5).
\setlength{\unitlength}{1pt}
\begin{picture}(135,90)(-25,-10)
  \put(0,35){\circle*{6}} \put(40,70){\circle*{6}} \put(40,0){\circle*{6}}
  \put(80,70){\circle*{6}} \put(80,0){\circle*{6}}
  \put(0,35){\line(4,3){40}}   % v1 v2
  \put(0,35){\line(4,-3){40}}  % v1 v3
  \put(40,70){\line(0,-1){70}} % v2 v3
  \put(40,70){\line(1,0){40}}  % v2 v4
  \put(40,0){\line(2,3){40}}   % v3 v4
  \put(40,0){\line(1,0){40}}   % v3 v5
  \put(80,70){\line(0,-1){70}} % v4 v5
  \put(-20,32){$v_1$} \put(33,76){$v_2$} \put(33,-12){$v_3$}
  \put(84,70){$v_4$} \put(84,-3){$v_5$}
\end{picture}

\caption{The graph $G_{\mathrm{tb}}$ with the seven edges that the textbook's own
examples use: $v_1v_2$, $v_1v_3$, $v_2v_3$, $v_2v_4$, $v_3v_4$, $v_3v_5$, $v_4v_5$. It is a part
of the textbook's Fig.~1.2.5 (printed p.~22)~\cite{petrosyan2015grafneri}. The textbook's
figure also has the edge $v_1v_5$, which none of the textbook's examples uses and which this paper does not use.}
\label{fig:textbook}
\end{figure}

\subsection{Graphs and the adjacency matrix}\label{sec:bg-graphs}

\begin{plainwords}
A village is a list of houses and a list of roads. Each road joins two
different houses. Two roads never join the same pair. Roads have no
direction. The list of houses is finite and has at least one house.
\end{plainwords}

\begin{definition}[Simple undirected graph; textbook Definition~1.1.1, p.~9]\label{def:graph}
A \emph{simple undirected graph} is a pair $G=(V,E)$, where $V$ is a finite nonempty set,
the \emph{vertices}, and $E$ is a set of unordered pairs of distinct vertices, the
\emph{edges}. The edge $\{u,v\}$ is written $uv$~\cite[Section~1.1, Definition~1.1.1, p.~9]{petrosyan2015grafneri}.
\end{definition}

\begin{remark}\label{rem:graph}
Because the pair must consist of distinct vertices, there is no loop (an edge from a
vertex to itself). Because $E$ is a set, no pair occurs twice, so there is no repeated edge.
Because the pair is unordered, $uv$ and $vu$ are the same edge. Isolated
vertices (vertices with no edge) are allowed, since $E$ may be any set of such pairs.
An edge is fixed by its two ends. This is why a walk can later be written by its vertices
alone (\Cref{rem:walk}): two different edges would need different pairs of ends. In this paper the word \emph{loop} always means such a self-edge.
\end{remark}

\begin{example}\label{ex:graph}
$V=\{n_1,n_2,n_3\}$ and $E=\{n_1n_2,\,n_2n_3\}$. This is the graph $B$ of the board matrix, three vertices in a row.
\end{example}

\begin{nonexample}\label{nex:graph}
Two vertices joined by two edges, or an edge from one vertex back to itself. The
first repeats an edge, the second is a loop. Neither is a simple graph.
\end{nonexample}

\begin{plainwords}
Two houses are neighbours when one road joins them directly. A house is never
its own neighbour. Being neighbours says nothing about longer trips.
\end{plainwords}

\begin{definition}[Adjacent; textbook Definition~1.1.4, p.~11]\label{def:adjacent}
Vertices $u$ and $v$ are \emph{adjacent} if $uv\in E$~\cite[Definition~1.1.4, p.~11]{petrosyan2015grafneri}.
Then each is called a \emph{neighbour} of the other.
\end{definition}

\begin{remark}\label{rem:adjacent}
Adjacency is symmetric ($uv=vu$), and no vertex is adjacent to itself, since there are no loops. It says nothing about longer connections.
\end{remark}

\begin{example}\label{ex:adjacent}
In $B$, the vertices $n_1$ and $n_2$ are adjacent.
\end{example}

\begin{nonexample}\label{nex:adjacent}
In $B$, the vertices $n_1$ and $n_3$ are not adjacent, although a walk of length $2$ joins them.
\end{nonexample}

\begin{plainwords}
Make a square table with one row and one column for every house. Write a one
where the row house and the column house share a road. Write a zero everywhere
else. This table describes the whole village.
\end{plainwords}

\begin{definition}[Adjacency matrix; textbook p.~11, after Definition~1.1.6; board, first line]\label{def:adjmatrix}
A \emph{matrix} is a rectangular table of numbers; its entry in row $i$ and
column $j$ is written with the indices $i,j$. The \emph{adjacency matrix} of $G$ is the $N\times N$ matrix $A$ with
$A_{ij}=1$ if $n_i$ and $n_j$ are adjacent, and $A_{ij}=0$ otherwise~\cite[p.~11]{petrosyan2015grafneri}.
The board writes the same rule as $A_{ij}=1$ if $e_{ij}\in V$ and $0$ if $e_{ij}\notin V$, with its $V$ the edge set~\cite{whiteboard1,whiteboard2}.
\end{definition}

\begin{remark}\label{rem:adjmatrix}
The definition allows only the entries $0$ and $1$. It depends on how the vertices are numbered:
renumbering permutes rows and columns together. \kindtag{derived here}
\end{remark}

\begin{example}\label{ex:adjmatrix}
The board matrix, with rows $(0\,1\,0)$, $(1\,0\,1)$, $(0\,1\,0)$:
\[
A=\MBoardA,
\]
is the adjacency matrix of $B$, since $A_{12}=A_{21}=1$, $A_{23}=A_{32}=1$ and all other entries are $0$.
\end{example}

\begin{nonexample}\label{nex:adjmatrix}
The matrix with rows $(0\,1\,0)$, $(0\,0\,1)$, $(0\,1\,0)$ is not the adjacency matrix of any
simple undirected graph: its entry $A_{12}$ is $1$ but $A_{21}$ is $0$, which would
need a one-way edge. A matrix with a $1$ on the diagonal (a loop) or an entry $2$ (a double edge) is not one either.
\end{nonexample}

\begin{plainwords}
The table has some easy features. It looks the same when flipped over its
diagonal. The diagonal holds only zeros. And any table of ones and zeros with
these two features comes from some village.
\end{plainwords}

\begin{proposition}[Properties of $A$]\label{prop:Aproperties}
For the adjacency matrix $A$ of a simple undirected graph:
\begin{enumerate}
\item every entry is $0$ or $1$;
\item $A$ is symmetric: $A_{ij}=A_{ji}$ for all $i,j$;
\item the diagonal is zero: $A_{ii}=0$ for all $i$;
\item conversely, every $N\times N$ matrix of zeros and ones that is symmetric with zero diagonal is the adjacency matrix of some simple graph with numbered vertices.
\end{enumerate}
These are stated in the textbook~\cite[pp.~11--12]{petrosyan2015grafneri}.
\end{proposition}

\begin{proof}
Items 1 to 3 are read off \Cref{def:adjmatrix}: adjacency is symmetric and irreflexive (\Cref{rem:adjacent}).
Item 4: join $n_i$ and $n_j$ exactly when the entry is $1$. Symmetry makes this a legitimate set of unordered pairs, and the zero diagonal excludes loops.
\end{proof}

\begin{example}\label{ex:Aproperties}
The board matrix is symmetric with zero diagonal, as items 2 and 3 require.
\end{example}

\begin{nonexample}\label{nex:Aproperties}
The matrix with rows $(0\,1\,0)$, $(1\,0\,1)$, $(1\,1\,0)$ is not symmetric (its entry $A_{31}$ is $1$ and $A_{13}$ is $0$), so it is not an adjacency matrix.
\end{nonexample}

\begin{plainwords}
The busy-ness of a house is the number of roads that meet there. We call that
number the degree of the house.
\end{plainwords}

\begin{definition}[Degree; textbook Definition~1.2.1, p.~16]\label{def:degree}
The \emph{degree} of a vertex $v$, written $d_G(v)$ or $d(v)$ in the textbook, is the number of edges incident to $v$, that is, the number of edges that contain $v$~\cite[Definition~1.2.1, p.~16]{petrosyan2015grafneri}. This paper writes $\deg(n_i)$, to keep $d$ for distance (\Cref{def:distance}).
\end{definition}

\begin{remark}\label{rem:degree}
In a simple graph the degree equals the number of neighbours, since each neighbour
is joined to $n_i$ by exactly one edge. \kindtag{derived here} The textbook places the definition of degree after the adjacency matrix, in Section~1.2, before the walks; this paper keeps that order.
\end{remark}

\begin{example}\label{ex:degree}
In $B$: $\deg(n_1)=1$, $\deg(n_2)=2$, $\deg(n_3)=1$.
\end{example}

\begin{nonexample}\label{nex:degree}
$\deg(n_1)$ is not $2$ in $B$, although two other vertices can be reached from $n_1$. Only edges at $n_1$ count, and $n_3$ is reached by passing through $n_2$.
\end{nonexample}

\begin{plainwords}
Adding up one row of the table counts the roads at that house. Adding up the whole
table counts every road twice, once from each end.
\end{plainwords}

\begin{proposition}[Row sums and the total]\label{prop:rowsums}
For the adjacency matrix $A$ of a simple undirected graph: (a)~the sum of row $i$, and of column $i$, equals $\deg(n_i)$; (b)~the sum of all entries equals $2|E|$. Part (b) is the textbook's Theorem~1.2.1 of Euler, $\sum_v d_G(v)=2|E|$~\cite[Theorem~1.2.1, p.~16]{petrosyan2015grafneri}, read through (a). Part (a) is stated, for graphs that may have loops, by Paramadevan and Sotheeswaran~\cite{paramadevan2021adjacency}. The proof here is \kindtag{derived here}.
\end{proposition}

\begin{proof}
(a) The sum of row $i$ is the number of $j$ with $A_{ij}=1$, that is, the number of neighbours of $n_i$. By \cref{rem:degree} this is $\deg(n_i)$. Column $i$ gives the same by symmetry (\cref{prop:Aproperties}). (b) Add (a) over all $i$. Each edge $n_in_j$ is counted once in $\deg(n_i)$ and once in $\deg(n_j)$, so twice in all.
\end{proof}

\begin{example}\label{ex:rowsums}
For the board matrix the row sums are $1,2,1$, the degrees of \cref{ex:degree}, and the sum of all entries is $4=2|E|$ with $|E|=2$.
\end{example}

\begin{nonexample}\label{nex:rowsums}
The matrix of \cref{nex:Aproperties} has row sums $1,2,2$, but it is not an adjacency matrix, so part (a) does not apply to it.
\end{nonexample}

\subsection{Walks, trails, paths and cycles}\label{sec:bg-walks}

The textbook uses the word ``path'' for a walk with no repeated edge
(Definition~1.2.12) and ``simple path'' for a walk with no repeated vertex
(Definition~1.2.13)~\cite{petrosyan2015grafneri}. English texts and the course
convention use ``path'' for the second~\cite{grinberg2025graph}. This paper uses
\emph{trail} for the first and \emph{path} for the second, as announced in
\cref{sec:conventions}. \Cref{tab:conventions} in \cref{sec:answer3} shows the
three conventions side by side.

\begin{plainwords}
A trip is a list of houses. Each house on the list shares a road with the one
before it. You may pass a house twice. You may drive a road twice. The length
of a trip is the number of roads driven. A road driven twice counts twice.
Staying at home is a trip of length zero. Books call a trip like this a walk.
\end{plainwords}

\begin{definition}[Walk; textbook Definition~1.2.10, p.~21]\label{def:walk}
A \emph{walk of length $k$ from $u_0$ to $u_k$}, also called a \emph{$(u_0,u_k)$-walk}, is a
sequence $u_0,e_1,u_1,\dots,u_{k-1},e_k,u_k$ of vertices and edges with $e_i=u_{i-1}u_i$ for $1\le i\le k$~\cite[Definition~1.2.10, p.~21]{petrosyan2015grafneri}.
Its \emph{length} is $k$, the number of edges, each edge counted as many times as it occurs~\cite[remark after Definition~1.2.10, p.~22]{petrosyan2015grafneri}.
\end{definition}

\begin{remark}\label{rem:walk}
Vertices and edges may repeat in a walk, and the length counts every use of an edge~\cite[p.~22]{petrosyan2015grafneri}. The only thing forbidden is a jump:
two consecutive vertices must be adjacent. In a simple graph the edge $e_i$ is fixed by its two ends (\Cref{rem:graph}), so a walk can be written by its vertices alone, $u_0,u_1,\dots,u_k$, and two walks are different exactly when these lists differ at some position. \kindtag{derived here}
The textbook does not say in words that $k=0$ is allowed. This paper allows it, because the textbook's own property $d(u,u)=0$ (\Cref{def:distance}) needs a walk of length $0$. \kindtag{interpretation, qualified}
\end{remark}

\begin{example}\label{ex:walk}
In $B$ the list $1,2,1,2,3$ is a walk from $n_1$ to $n_3$ of length $4$: its steps $12$, $21$, $12$, $23$ are all edges. Vertices $1$ and $2$ occur twice, and the edge $12$ is used twice. The textbook's example is $v_1,v_3,v_2,v_4,v_3,v_2,v_4,v_3,v_5$ in $G_{\mathrm{tb}}$, a walk of length $8$~\cite[p.~22]{petrosyan2015grafneri}.
\end{example}

\begin{nonexample}\label{nex:walk}
In $B$ the list $1,3$ is not a walk, because $13$ is not an edge of $B$ (the matrix has $A_{13}=0$). It would be a walk of the drawing $T$. The textbook's non-example is $v_3,v_1,v_4,v_2,v_3,v_5$, which is not a walk of $G_{\mathrm{tb}}$~\cite[p.~22]{petrosyan2015grafneri}.
\end{nonexample}

\begin{plainwords}
A trip is closed when it ends at the house where it began. That is the only new
rule. The trip may still pass houses and roads many times.
\end{plainwords}

\begin{definition}[Closed walk; textbook Definition~1.2.11, p.~22]\label{def:closedwalk}
A $(u_0,u_k)$-walk is \emph{closed} if $u_0=u_k$~\cite[Definition~1.2.11, p.~22]{petrosyan2015grafneri}.
\end{definition}

\begin{remark}\label{rem:closedwalk}
A closed walk allows everything a walk allows; it forbids only a different start and end. A walk of length $0$ is closed. In a simple graph a closed walk cannot have length $1$ (that would be a loop), and one of length $2$ must be $u,w,u$ and so uses one edge twice. \kindtag{derived here}
\end{remark}

\begin{example}\label{ex:closedwalk}
In $B$ the walks $2,1,2$ (length $2$) and $2,3,2,1,2$ (length $4$) are closed.
\end{example}

\begin{nonexample}\label{nex:closedwalk}
In $B$ the walk $1,2,3$ is not closed: it starts at $1$ and ends at $3$.
\end{nonexample}

\begin{plainwords}
Now add one rule. Never drive the same road twice. You may still pass a house
twice, if you return by another road. English books call such a trip a trail.
Our textbook calls it a path. This paper says trail, so that the word path can
keep one meaning.
\end{plainwords}

\begin{definition}[Trail; textbook Definition~1.2.12, p.~22, where it is called a ``path'']\label{def:trail}
A $(u_0,u_k)$-walk is a \emph{trail} if its edges $u_0u_1,\dots,u_{k-1}u_k$ are pairwise distinct. The textbook calls this a \emph{path}~\cite[Definition~1.2.12, p.~22]{petrosyan2015grafneri}; English usage calls it a trail~\cite{grinberg2025graph}.
\end{definition}

\begin{remark}\label{rem:trail}
A trail allows repeated vertices and forbids repeated edges. A vertex can recur in a trail only if the walk left it by one edge and came back by another, so the trail contains a closed trail, a detour that comes back to the same vertex. In $B$ every trail is also a path in the sense of \cref{def:path}: a repeated vertex would need a closed trail, which needs at least three edges (\cref{rem:closedwalk}), and $B$ has two. So $B$ cannot show the difference; $H$ and the textbook's graph can. \kindtag{derived here}\label{rem:Bsimple}
\end{remark}

\begin{example}\label{ex:trail}
In $G_{\mathrm{tb}}$ the list $v_1,v_3,v_2,v_4,v_3,v_5$ is a trail: its edges $v_1v_3$, $v_3v_2$, $v_2v_4$, $v_4v_3$, $v_3v_5$ are distinct, although $v_3$ occurs twice~\cite[p.~22]{petrosyan2015grafneri}. In $H$ the list $4,3,1,2,3$ is a trail: its edges $43$, $31$, $12$, $23$ are distinct, although vertex $3$ occurs twice.
\end{example}

\begin{nonexample}\label{nex:trail}
The walk $v_1,v_3,v_2,v_4,v_3,v_2,v_4,v_3,v_5$ of length $8$ is not a trail: the edge $v_3v_2$ occurs twice~\cite[p.~22]{petrosyan2015grafneri}. In $B$ the walk $1,2,1,2,3$ is not a trail, since the edge $12$ occurs twice; in $H$ the walk $1,2,3,2$ is not a trail, since the edge $23$ occurs twice.
\end{nonexample}

\begin{plainwords}
Now add a stronger rule. Never pass the same house twice. Then you never drive
a road twice either. A trip like this is a route. It goes straight from start to end. It has
no detour in it. Our course calls it a path. The textbook calls it a simple path.
\end{plainwords}

\begin{definition}[Path; textbook Definition~1.2.13, p.~22, where it is called a ``simple path'']\label{def:path}
A $(u_0,u_k)$-walk is a \emph{path} if its vertices $u_0,\dots,u_k$ are pairwise distinct. The textbook calls this a \emph{simple path} and writes $P_n$ for the simple path on $n$ vertices (here $n$ counts the vertices of that path, not of the whole graph)~\cite[Definition~1.2.13, p.~22]{petrosyan2015grafneri}. Its length is $k$.
\end{definition}

\begin{remark}\label{rem:path}
A path forbids repeated vertices, hence also repeated edges (\Cref{prop:pathtrail}). It allows length $0$ (a single vertex is a path from that vertex to itself) and length $1$ (one edge). It does not allow a closed path of positive length, because the first and last vertex would be equal.
The textbook defines a simple path as a trail with distinct vertices; by \cref{prop:pathtrail} this is the same as a walk with distinct vertices, so the two ways of saying it agree. That the \emph{course} uses ``path'' in this sense is \kindtag{interpretation, qualified}: it rests on the user's report of the assignment text, and the board does not say.
\end{remark}

\begin{example}\label{ex:path}
In $B$ the list $1,2,3$ is a path from $n_1$ to $n_3$ of length $2$. In $H$ the list $4,3,1,2$ is a path of length $3$. The textbook's example is $v_1,v_3,v_4,v_5$ in $G_{\mathrm{tb}}$, a simple path~\cite[p.~22]{petrosyan2015grafneri}.
\end{example}

\begin{nonexample}\label{nex:path}
In $B$ the list $1,2,3,2$ is a walk but not a path: vertex $2$ occurs twice. In $G_{\mathrm{tb}}$ the list $v_1,v_3,v_2,v_4,v_3,v_5$ is a trail but not a simple path, since $v_3$ occurs twice~\cite[p.~22]{petrosyan2015grafneri}. In $H$ the trail $4,3,1,2,3$ is not a path for the same reason.
\end{nonexample}

\begin{plainwords}
A trip that never passes a house twice never drives a road twice. Here is why.
Take two roads on the trip, an earlier one and a later one. The later road ends
at a house that the earlier road does not touch. So the two roads are different.
\end{plainwords}

\begin{proposition}[A path is a trail]\label{prop:pathtrail}
In a simple graph, a walk with pairwise distinct vertices has pairwise distinct edges. So every path is a trail. \kindtag{derived here}
\end{proposition}

\begin{proof}
Let $u_0,\dots,u_k$ have pairwise distinct vertices, and take two steps with indices $a<b$, that is, the edges $u_{a-1}u_a$ and $u_{b-1}u_b$. The vertex $u_b$ has an index larger than the indices $a-1$ and $a$, so it differs from both $u_{a-1}$ and $u_a$. But $u_b$ is an end of the second edge and is not an end of the first. So the two edges are different. (Grinberg leaves this statement as an exercise~\cite{grinberg2025graph}.)
\end{proof}

\begin{plainwords}
A cycle is a closed trip that never drives a road twice. It may pass a house
twice on the way. An English book calls it a closed trail.
\end{plainwords}

\begin{definition}[Cycle; textbook Definition~1.2.14, p.~22]\label{def:cycle}
A trail that is a closed walk is called a \emph{closed path} or \emph{cycle}~\cite[Definition~1.2.14, p.~22]{petrosyan2015grafneri}. In the words of this paper, a \emph{cycle} is a closed trail: a $(u_0,u_k)$-walk with $u_0=u_k$ and pairwise distinct edges. Vertices may repeat.
\end{definition}

\begin{remark}\label{rem:cycle}
The definition allows repeated vertices (a figure eight is a cycle) and forbids repeated edges and open walks. English usage differs: Grinberg's ``cycle'' needs distinct vertices and is the textbook's simple cycle~\cite{grinberg2025graph}. The textbook is silent on whether a closed trail of length $0$ is a cycle; this paper does not count it \kindtag{interpretation, qualified}.
\end{remark}

\begin{example}\label{ex:cycle}
In $G_{\mathrm{tb}}$ the list $v_3,v_2,v_1,v_3,v_5,v_4,v_3$ is closed and has six distinct edges, so it is a cycle in this sense; it is not a simple cycle, since $v_3$ occurs three times~\cite[p.~22]{petrosyan2015grafneri}.
\end{example}

\begin{nonexample}\label{nex:cycle}
In $B$ the closed walk $1,2,1$ is not a cycle, because the edge $12$ is used twice; in $H$ the closed walk $4,3,4$ is not a cycle for the same reason. In fact $B$ has no cycle at all: a cycle is a closed trail, in a simple graph a closed trail of positive length needs at least three distinct edges (\cref{rem:closedwalk}), and $B$ has only two edges. \kindtag{derived here}
\end{nonexample}

\begin{plainwords}
A simple cycle is stricter. It passes no house twice, except the start, which is
also the end. It is a ring of houses, each joined by a road to the next.
\end{plainwords}

\begin{definition}[Simple cycle; textbook Definition~1.2.15, p.~22]\label{def:simplecycle}
A cycle is \emph{simple} if only its first and last vertices coincide. The simple cycle on $n$ vertices is written $C_n$, with $n\ge3$ (again $n$ counts the vertices of the cycle)~\cite[Definition~1.2.15, p.~22]{petrosyan2015grafneri}.
\end{definition}

\begin{remark}\label{rem:simplecycle}
A simple cycle forbids any repeated vertex other than the closing one. The condition $n\ge3$ is forced in a simple graph: a closed trail cannot have length $1$ or $2$ (\Cref{rem:closedwalk}). \kindtag{derived here}
\end{remark}

\begin{example}\label{ex:simplecycle}
In $G_{\mathrm{tb}}$ the list $v_1,v_3,v_4,v_2,v_1$ is a simple cycle~\cite[p.~22]{petrosyan2015grafneri}. In the board drawing $T$ and in $H$ the list $1,2,3,1$ is a simple cycle of length $3$ (by script).
\end{example}

\begin{nonexample}\label{nex:simplecycle}
The cycle $v_3,v_2,v_1,v_3,v_5,v_4,v_3$ of \cref{ex:cycle} is not simple, because $v_3$ occurs three times~\cite[p.~22]{petrosyan2015grafneri}.
\end{nonexample}

\begin{plainwords}
The distance between two houses is the number of roads on the shortest trip
between them. If no trip exists, the distance is called infinite. A house is at
distance zero from itself. The distance there equals the distance back, because
any trip can be driven backwards. Going by way of a third house is never shorter
than the shortest trip.
\end{plainwords}

\begin{definition}[Distance; textbook Definition~1.2.16, pp.~22--23]\label{def:distance}
The \emph{distance} $d(u,v)=d_G(u,v)$ between vertices $u$ and $v$ is the length of a shortest $(u,v)$-trail if at least one $(u,v)$-trail exists, and $+\infty$ otherwise~\cite[Definition~1.2.16, pp.~22--23]{petrosyan2015grafneri}. (The textbook's word is ``path'', its Definition~1.2.12.)
\end{definition}

\begin{remark}\label{rem:distance}
The textbook lists three properties~\cite[p.~23]{petrosyan2015grafneri}:
\begin{enumerate}
\item $d(u,v)\ge0$, and $d(u,v)=0$ if and only if $u=v$;
\item $d(u,v)=d(v,u)$;
\item $d(u,v)\le d(u,w)+d(w,v)$.
\end{enumerate}
So if any two vertices are joined by a trail, $(V,d)$ is a metric space~\cite{petrosyan2015grafneri}; a \emph{metric space} is just a set with a distance that satisfies these three properties, and the word is not used again. The definition allows the value $+\infty$, which is not an integer; this is how a missing trail shows up. It forbids negative values and any value that is not the length of an actual trail. The shortest trail, the shortest path and the shortest walk between two vertices have the same length, because a shortest walk contains a path that is no longer than itself (\cref{lem:walkpath}); so the definition does not depend on whether ``path'' is read as trail or path (proof in \cref{app:distance}). \kindtag{derived here}
\end{remark}

\begin{example}\label{ex:distance}
In $B$: $d(1,2)=1$, $d(2,3)=1$, $d(1,3)=2$, and $d(i,i)=0$ for each $i$ (by script). The walk $1,2,1,2,3$ has length $4$, but $d(1,3)$ is $2$, the shortest.
\end{example}

\begin{nonexample}\label{nex:distance}
The number $4$ is not the distance from $1$ to $3$ in $B$, even though a walk of length $4$ exists: distance is the minimum, not any length. In the graph $D$ with vertices $1,2,3$ and the single edge $12$ (\cref{fig:oneedge}), $d(1,3)$ is not a number at all; it is $+\infty$ (by script).
\end{nonexample}

\begin{figure}[ht]
\centering
\setlength{\unitlength}{1pt}
\begin{picture}(130,40)(0,-5)
  \put(10,15){\circle*{6}} \put(65,15){\circle*{6}} \put(120,15){\circle*{6}}
  \put(10,15){\line(1,0){55}}
  \put(7,22){$1$} \put(62,22){$2$} \put(117,22){$3$}
\end{picture}

\caption{The graph $D$: three vertices and the single edge $12$. Vertex $3$ is isolated. It is the board graph $B$ with the edge $23$ removed.}
\label{fig:oneedge}
\end{figure}

\subsection{From walks to paths}\label{sec:bg-lemma}

\begin{plainwords}
Suppose you can drive from one house to another. Then you can do so without
passing any house twice. Look at your trip. If it passes some house twice, you
went on a detour that came back to the same house. Cut the detour out. What is left is still a trip between the
same two houses. It is shorter. Keep cutting detours until no house repeats. The trip
that remains is never longer than the first one.
\end{plainwords}

\begin{lemma}[Every walk contains a path; textbook Lemma~1.2.1, p.~23, for $u\ne v$]\label{lem:walkpath}
Let $u$ and $v$ be vertices of a simple graph, equal or different. If there is a $(u,v)$-walk of length $k$, then there is a $(u,v)$-path of length at most $k$. All vertices and edges of this path are vertices and edges of the walk.
\end{lemma}

The textbook states the lemma for $u\ne v$, as ``from any $(u,v)$-walk one can extract a simple $(u,v)$-path'', without a length bound, and proves it by induction on the length of the walk~\cite[Lemma~1.2.1, p.~23]{petrosyan2015grafneri}. Grinberg's Corollary~2.9.7 covers $u=v$ and gives the length bound~\cite[Corollary~2.9.7, p.~39]{grinberg2025graph}. The case $u=v$ is needed here, because a closed walk such as $2,1,2$ in $B$ has $u=v$; the path it contains is the single vertex, of length $0$. The proof below is the paper's own \kindtag{derived here}; it follows the textbook's method.

\begin{proof}
First an example of the idea. In $B$ the walk $1,2,1,2,3$ has the vertex $1$ at positions $0$ and $2$; removing the part between them leaves $1,2,3$, a shorter walk with the same ends, and it is a path.

We prove by strong induction on $k$ the statement $P(k)$: every $(u,v)$-walk of length $k$, for any vertices $u$ and $v$, contains a $(u,v)$-path of length at most $k$.

\emph{Base case, $k=0$.} A walk of length $0$ is a single vertex $u$, so $v=u$. A single vertex has no repeated vertex, so the walk is itself a path, of length $0\le0$.

\emph{Induction hypothesis.} Let $k\ge1$ and assume $P(m)$ holds for every $m<k$.

\emph{Step.} Let $u_0,u_1,\dots,u_k$ be a $(u,v)$-walk, so $u_0=u$ and $u_k=v$.
Case~1: the vertices $u_0,\dots,u_k$ are pairwise distinct. Then the walk is a path of length $k\le k$.
Case~2: some vertex occurs twice, so $u_i=u_j$ for some $i<j$. Form the list $W'=u_0,\dots,u_i,u_{j+1},\dots,u_k$. If $j=k$, the part after $u_i$ is empty and $W'$ ends at $u_i$. We check that $W'$ is a $(u,v)$-walk. (a)~It starts at $u_0=u$. (b)~It ends at $v$: if $j<k$ its last vertex is $u_k=v$; if $j=k$ its last vertex is $u_i=u_j=u_k=v$. (c)~Consecutive vertices are adjacent: before the cut and after the cut they are consecutive vertices of the old walk, and at the join the pair is $(u_i,u_{j+1})=(u_j,u_{j+1})$, which is a step of the old walk. The length of $W'$ is $i+(k-j)$, and since $i<j$ this is less than $k$. By the induction hypothesis $W'$ contains a $(u,v)$-path of length at most that, hence at most $k$. Every vertex and edge of that path is one of $W'$, and so one of the old walk.
\end{proof}

Why the proof is valid: each cut removes at least one step, since $j-i\ge1$, so the length drops by at least $1$ and the process stops. A walk that cannot be cut has no repeated vertex, which is the definition of a path.

\begin{remark}\label{rem:lemmaexample}
The textbook's walk of length $8$ in $G_{\mathrm{tb}}$ is $1,3,2,4,3,2,4,3,5$, written by the indices of its vertices, with positions counted from $0$. Positions $1$ and $4$ both hold vertex $3$ (so $u_1=u_4=3$). The cut removes the detour in brackets and keeps the rest:
\[
1,\,3,\,[2,\,4,\,3],\,2,\,4,\,3,\,5 \;\longrightarrow\; 1,\,3,\,2,\,4,\,3,\,5 \quad(\text{length }5).
\]
In the shorter walk positions $1$ and $4$ again both hold vertex $3$, and the cut leaves a path:
\[
1,\,3,\,[2,\,4,\,3],\,5 \;\longrightarrow\; 1,\,3,\,5 \quad(\text{a path of length }2\le8).
\]
In $B$ the walk $1,2,1,2,3$ cuts to $1,[2,1],2,3 \to 1,2,3$; the closed walk $2,1,2$ cuts to the single vertex $2$. (Checked by script, \cref{app:scripts}.)
\end{remark}

\begin{plainwords}
A trip that never repeats a house can pass each house at most once. So when the
number of houses is fixed, such a trip is short. It uses at most one road fewer
than the number of houses.
\end{plainwords}

\begin{corollary}[Length bound]\label{cor:bound}
In a graph with $N$ vertices every path has at most $N-1$ edges. Hence, if there is a walk from $u$ to $v$, there is a walk from $u$ to $v$ of length at most $N-1$. \kindtag{derived here}
\end{corollary}

\begin{proof}
A path of length $k$ lists $k+1$ pairwise distinct vertices, and there are only $N$ vertices, so $k+1\le N$, that is $k\le N-1$. For the second part, \cref{lem:walkpath} gives a path inside the walk, of length at most that of the walk and at most $N-1$; a path is itself a walk. Grinberg gives the same argument~\cite[Section~2.9.4, p.~40]{grinberg2025graph}.
\end{proof}

For $B$, $N=3$, so a path has at most $2$ edges; the longest path in $B$ has $2$ edges (by script).

\subsection{Connected graphs and components}\label{sec:bg-connected}

\begin{plainwords}
Sometimes you can reach every house from every other house. Then the village is
called connected. If one house is cut off, the village is not connected.
\end{plainwords}

\begin{definition}[Connected; textbook Definition~2.1.1, p.~48; board, last line]\label{def:connected}
A graph $G$ is \emph{connected} if for any two vertices $u$ and $v$ of $G$ there is a $(u,v)$-path in $G$~\cite[Definition~2.1.1, p.~48]{petrosyan2015grafneri}. The board says the same: ``Graph is called connected if for any $n_i,n_j$ there exists a path between nodes''~\cite{whiteboard1,whiteboard2}.
\end{definition}

\begin{remark}\label{rem:connected}
In the textbook the word ``path'' here is Definition~1.2.12, a trail; the board does not say which meaning it uses. The next proposition shows that the choice does not change which graphs are connected. ``Any two vertices'' includes $u=v$; for $u=v$ the one-vertex walk is a $(u,u)$-path, so these pairs are always satisfied, and reading the definition as ``any two different vertices'' changes nothing. \kindtag{interpretation; the textbook does not discuss $u=v$}
The definition forbids exactly one thing: a pair of vertices with no path between them.
\end{remark}

\begin{example}\label{ex:connected}
The board graph $B$ is connected: the pair $1,3$ is joined by the path $1,2,3$, and the pairs $1,2$ and $2,3$ by one edge each. The textbook's examples of connected graphs are the complete graphs, the complete bipartite graphs, and the three graphs of its Fig.~2.1.1~\cite[p.~48]{petrosyan2015grafneri}.
\end{example}

\begin{nonexample}\label{nex:connected}
The graph $D$ (\cref{fig:oneedge}) is not connected: no walk starts at vertex $3$ and ends elsewhere, so the pair $1,3$ has no path. The textbook's non-example is the graph of its Fig.~2.1.2, which has four components~\cite[p.~49]{petrosyan2015grafneri}.
\end{nonexample}

\begin{plainwords}
It does not matter which kind of trip we ask for. If any trip links two houses,
a route links them. Every route is a trail. Every trail is a trip. So all three
questions have the same answer.
\end{plainwords}

\begin{proposition}[Three readings of ``connected'' agree]\label{prop:readings}
For vertices $u,v$ of a simple graph the following are equivalent: (a) there is a $(u,v)$-walk; (b) there is a $(u,v)$-trail; (c) there is a $(u,v)$-path. Hence $G$ is connected in the textbook's sense (trails) if and only if every two vertices are joined by a path, if and only if every two vertices are joined by a walk. \kindtag{derived here from \cref{lem:walkpath}}
\end{proposition}

\begin{proof}
(c)$\Rightarrow$(b): every path is a trail (\cref{prop:pathtrail}). (b)$\Rightarrow$(a): every trail is a walk, by definition. (a)$\Rightarrow$(c): \cref{lem:walkpath}. Grinberg states the equivalence of (a) and (c) as his Proposition~2.9.10~\cite[p.~42]{grinberg2025graph}; the lemma behind (a)$\Rightarrow$(c) is the textbook's Lemma~1.2.1 for $u\ne v$ and Grinberg's Corollary~2.9.7 in general.
\end{proof}

\begin{plainwords}
A village that is not connected falls into pieces. Inside one piece every house
can reach every other house. No road joins two different pieces. Each piece is
called a component. A connected village has exactly one piece.
\end{plainwords}

\begin{definition}[Connected component; textbook, p.~49, paragraph before Remark~2.1.1]\label{def:component}
On the vertex set $V$ write $u\sim v$ if there is a $(u,v)$-path in $G$. This relation is \emph{reflexive} (every vertex is related to itself), \emph{symmetric} (if $u\sim v$ then $v\sim u$) and \emph{transitive} (if $u\sim v$ and $v\sim w$ then $u\sim w$)~\cite[p.~49]{petrosyan2015grafneri}. Here is why, reading ``path'' as walk, which is allowed by \cref{prop:readings}: the one-vertex walk gives reflexivity, reversing a walk gives symmetry, and joining a $(u,v)$-walk to a $(v,w)$-walk gives a $(u,w)$-walk, which gives transitivity \kindtag{derived here}. Hence $V$ splits into \emph{classes} $V_1,\dots,V_p$: groups of vertices such that each vertex lies in exactly one group, and two vertices lie in the same group exactly when they are related. The subgraphs $G_j=G[V_j]$ (the vertices of $V_j$ with every edge of $G$ that has both ends in $V_j$) are the \emph{connected components} of $G$~\cite[p.~49]{petrosyan2015grafneri}. The textbook gives this in prose, without a definition number.
\end{definition}

\begin{remark}\label{rem:component}
A component is a whole class: a vertex reachable from a vertex of $V_j$ lies in $V_j$. The textbook's Remark~2.1.1 states that $G$ is connected if and only if it has exactly one connectivity component~\cite[Remark~2.1.1, p.~49]{petrosyan2015grafneri}. By \cref{prop:readings} the relation $\sim$ is the same whether ``path'' means trail, path or walk.
\end{remark}

\begin{example}\label{ex:component}
The graph $D$ has the two components $G[\{1,2\}]$ and $G[\{3\}]$; the board graph $B$ has the single component $G[\{1,2,3\}]$ (checked by breadth-first search in the script, \cref{app:scripts}).
\end{example}

\begin{nonexample}\label{nex:component}
$G[\{1\}]$ is not a component of $D$, although it is a connected subgraph: vertex $2$ can be reached from vertex $1$, so $2$ lies in the same class as $1$. A component is a whole class, not any connected piece. \kindtag{derived here}
\end{nonexample}
